\subsection{Projektinitialisierung und Generierung}
\label{subsec:projektinitialisierung-generierung}

\texttt{init} erzeugt aus Profil und kompatiblen Capabilities ein eigenständiges
Produktprojekt. Die Auswahl kann geführt oder nicht interaktiv erfolgen;
Projektidentität, Ziel und Dry-Run sind konfigurierbar. Eine geänderte
Tauri-Identität erfordert einen gültigen Reverse-Domain-Identifier
\parencite{kleiveist2026profiles,kleiveist2026tooling}.

\begin{figure}[H]
  \centering
  \resizebox{0.72\textwidth}{!}{%
  \begin{tikzpicture}[node distance=5mm and 10mm]
    \node[casePrimary,minimum width=52mm] (master) {Master-Repository};
    \node[caseBox,minimum width=52mm,below=of master] (selection)
      {Profil + Capabilities\\optionale Produktidentität};
    \node[caseBox,minimum width=52mm,below=of selection] (features)
      {Abhängigkeiten auflösen\\Kombination validieren};
    \node[caseBox,minimum width=52mm,below=of features] (plan)
      {Scaffold-Plan\\Core + aktive Feature-Pfade};
    \node[caseBox,minimum width=52mm,below=of plan] (generation)
      {Pfade kopieren\\Konfiguration und Identität ableiten};
    \node[casePrimary,minimum width=52mm,below=of generation] (product)
      {eigenständiges Produktprojekt};
    \node[caseOptional,minimum width=42mm,right=of product] (manifest)
      {aktives Manifest\\\texttt{project-profile.toml}};

    \draw[caseFlow] (master) -- (selection);
    \draw[caseFlow] (selection) -- (features);
    \draw[caseFlow] (features) -- (plan);
    \draw[caseFlow] (plan) -- (generation);
    \draw[caseFlow] (generation) -- (product);
    \draw[caseSecondaryFlow] (product) -- (manifest);
  \end{tikzpicture}}
  \caption{Workflow der profilgesteuerten Projektgenerierung}
  \label{fig:project-generation-workflow}
  \smallskip
  {\footnotesize Quelle: Eigene Darstellung auf Grundlage von
  \textcite{kleiveist2026profiles} und \textcite{kleiveist2026tooling}.\par}
\end{figure}


Der Generator löst Abhängigkeiten auf, erstellt den Scaffold-Plan und prüft
Quellen und Ziel vor dem Schreiben. Er kopiert Core und aktive Features,
erzeugt Manifest, Frontend-Profil und \path{.env.example} und passt vorhandene
Produktmetadaten an. Transiente Verzeichnisse werden ausgelassen; Webprofile
verlieren Tauri-Skript und -CLI.
Bei angegebener Produktidentität werden nur tatsächlich vorhandene Frontend-,
Backend-, Compose-, Tauri- und Cargo-Metadaten angepasst. Der Umfang bleibt
damit vom gewählten Profil abhängig.

Im Master sind Ziele nur unter \path{.generated/} erlaubt; externe Ziele sind
möglich. Das Repository selbst und nichtleere Ziele werden abgewiesen,
\texttt{--dry-run} schreibt nichts. Eine Grenze bleibt: Nur der Dialog filtert
nach \texttt{selectable}; direktes \texttt{--with database} wird für
Backendprofile akzeptiert \parencite{kleiveist2026governance-snapshot}.
Diese Vorprüfungen halten Master und Ziel getrennt und verhindern teilweise
erzeugte Projekte bei bekannten Konflikten. Die Ableitung ist deterministisch
angelegt; der empirische Reproduzierbarkeitsnachweis bleibt dennoch auf den
geprüften Fall begrenzt.
\beginnerinput{workspace/statement/500_tooling_workflow/520}
